\section{Introduction}
\label{sec:introduction}

Consensus protocols enable mutually distrusting nodes to reach agreement in permissionless blockchain systems. Bitcoin’s Proof-of-Work (PoW) is the most prominent example, securing the network by requiring participants to solve computationally intensive puzzles and thereby increasing the cost of adversarial behavior. However, such systems have been widely criticized for performing large amounts of wasteful computation and for their environmental impact. For instance, the Bitcoin network alone is estimated to consume approximately 172 terawatt-hours of electricity annually, comparable to the total electricity consumption of a country such as Poland. These concerns have motivated the search for more efficient consensus mechanisms.

In this paper, we propose a novel consensus mechanism based on the increasing demand of machine learning (ML) workflows, which we term \emph{Proof-of-ML-Inference} (PoML). This consensus layer replaces the traditional PoW scheme used in blockchains with one grounded in useful computational effort. Rather than expending energy on hashing arbitrary nonces, nodes perform ML inference and generate zero-knowledge proofs (ZKPs) attesting to the correctness of these computations. As a result, the computational resources consumed by the network directly contribute to meaningful tasks, offering a practical realization of \emph{Proof-of-Useful-Work}.

Although prior proposals have explored machine learning as a form of useful work (e.g. 
\cite{pot,pocl,pouw}), they suffer from two fundamental limitations. First, their validation mechanisms are often inefficient: many require independent validators to recompute the original tasks in order to confirm correctness, incurring substantial overhead. Consequently, such protocols either rely on proof-of-stake mechanisms—where incorrect results are penalized through slashing—or delay block confirmation until validation is completed. Second, some proposals retain Bitcoin-like nonce grinding as an auxiliary mechanism, meaning a significant portion of mining effort remains wasteful. With recent advances in ZKP systems, however, it has become increasingly feasible to generate succinct and efficiently verifiable proofs even for large ML models. This enables inexpensive verification of correct computation without redundant execution or separate validators. We therefore propose to use ZKPs as the primary mechanism through which useful work contributes to consensus. \thomas{TODO improve the citation}

To fit the security requirements of a consensus protocol, PoML requires the target model class be stochastic in nature, so that nodes will never solve two instances of the same problem. It also needs to satisfy an important structural property:  \emph{Computational Independence}, which states that distinct random seeds produce execution traces whose intermediate computations cannot be reused across inferences. We show that \emph{diffusion models}---a popular text-to-image generation model class---is a natural and well-motivated instantiation of a model class that satisfies the above properties. In diffusion models, the stochastic sampling mechanism injects \emph{fresh Gaussian noise at every denoising step}, so that even when the same text prompt is reused, a new random seed produces intermediate activations that diverge from any prior execution. 

The main contributions of this work are as follows. 
\begin{enumerate}
    \item We propose PoML, a novel consensus mechanism in which block validity is determined entirely by the correctness of ML inference. Unlike prior useful-work proposals, PoML requires \emph{no nonce grinding} and \emph{no additional validation stage}---it operates using only standard Bitcoin primitives.
    \item We formally define the structural properties---uniform query cost and computational independence---that a model class must satisfy for secure PoML consensus, and show that stochastic denoising diffusion models satisfy these properties.
    \item We carry out a formal security reduction from PoML to the Bitcoin backbone protocol~\cite{GKL20}, proving that PoML inherits Bitcoin's core security guarantees---common prefix, chain quality, and chain growth---without modification.
\end{enumerate}
\thomas{Rik- Please check the contributions as above}
